\chapter{Dependency Map: Results, Assignments, and Research}

\section{What an arrow means in this map}

\begin{context}
A solid arrow means that the source result is used in proving or constructing the target
result. A dashed arrow marks an application or conceptual comparison: it does not claim
that the target result follows from the first. The assumptions written along each path are
part of the connection.
\end{context}

\section{Paths used in proofs}

\begin{dependencias}
\textbf{Continuity and inverses.}
\[
\boxed{\text{Neumann series}}
\longrightarrow
\boxed{\text{1.1(b): continuous inversion}}
\longrightarrow
\boxed{\text{Prop. 2.17: spectral permanence}}.
\]
The last arrow uses $(b+i/n)^{-1}\to b^{-1}$ in the ambient algebra and the fact that the
$C^*$-subalgebra is closed. Problem 1.1(a) supplies continuity of multiplication through
another route for Problems 1.3, 1.7, and 2.8.
\end{dependencias}

\begin{dependencias}
\textbf{Commutative reduction.}
\[
\begin{gathered}
\boxed{\text{Problem 1.2: commuting inverses}}
\quad+\quad
\boxed{\text{Problem 1.3: generated algebra}}\\[2mm]
\Downarrow\\[2mm]
\boxed{\text{Problem 1.9: inverse-closed subalgebra}}\\[2mm]
\Downarrow\\[2mm]
\boxed{\text{Prop. 2.13: real spectrum}}.
\end{gathered}
\]
The last proof also uses the $C^*$ identity, the fact that characters are real on
self-adjoint elements, and the description of the spectrum through characters.
Problem 1.9 supplies the commutative subalgebra that preserves the spectrum;
by itself it does not prove that the spectrum is real.
\end{dependencias}

\begin{dependencias}
\textbf{From functions to functional calculus.}
\[
\begin{gathered}
\boxed{\text{Problem 1.3(d): }C^*(a,1)\text{ commutative for normal $a$}}\\[2mm]
\quad+\quad
\boxed{\text{Prop. 2.17: spectral permanence}}
\quad+\quad
\boxed{\text{Thm. 2.20: Gelfand--Naimark}}\\[2mm]
\Downarrow\\[2mm]
\boxed{\text{Thm. 2.22: functional calculus}}.
\end{gathered}
\]
Problem 1.10 computes the spectrum of a continuous function and, together with functional
calculus and spectral permanence, gives Problem 2.4. Problems 2.2(a,c) give the forward
directions of Problem 2.6(c,b); the converses use normality and functional calculus.
\end{dependencias}

\begin{dependencias}
\textbf{Continuity as the operator varies.}
\[
\begin{gathered}
\boxed{\text{Problem 1.7: norm limit of normal elements is normal}}
\quad+\quad
\boxed{\text{Stone--Weierstrass}}\\[2mm]
\quad+\quad
\boxed{\text{isometric functional calculus}}\\[2mm]
\Downarrow\\[2mm]
\boxed{\text{Problem 2.8: }f(x_n)\to f(x)}.
\end{gathered}
\]
Convergence here is in norm. Problem 2.9 studies a different situation: a bounded,
increasing net of self-adjoint operators converges strongly, that is, vector by vector
in the norm of $H$.
\end{dependencias}

\section{Coverage of the exercises}

\small
\begin{longtable}{>{\bfseries}p{1.3cm}p{4.3cm}p{8.2cm}}
\toprule
Problem & Result proved & Exact use or scope of the connection\\
\midrule
\endhead
1.1 & Continuous multiplication and inversion & Multiplication supports Problems 1.3 and 1.7; inversion is used in Prop. 2.17, and multiplication again in Problem 2.8.\\
1.2 & Commutation with inverses and adjoints & Used when adjoining inverses in the construction for Problem 1.9.\\
1.3 & Generated algebra and commutation & For normal $a$, constructs commutative $C^*(a,1)$, used in Proposition 2.18 and Theorem 2.22.\\
1.4 & Quotient by a closed ideal & Makes $\mathcal B(X)/\K(X)$ a Banach algebra; the Fredholm--invertibility equivalence is the external Atkinson theorem.\\
1.5 & $C_0(X)$ closed in $C_b(X)$ & When $X$ is noncompact, it is nonunital; the general classification of $C_0(X)$ requires the nonunital version of the theorem in Section 2.\\
1.6 & $\spec(x^*)=\overline{\spec(x)}$ & Symmetry holds in involutive algebras; reality of a self-adjoint element's spectrum also needs the $C^*$ argument of Prop. 2.13.\\
1.7 & Normality is closed & The first assumption needed in Problem 2.8 to define $f(x)$ at the limit.\\
1.8 & Spectral inclusion under homomorphisms & Applied to $a^*a$, it enters the proof that $*$-homomorphisms are contractive (Thm. 2.8).\\
1.9 & Inverse-closed commutative subalgebra & Explicitly used in the proof of Prop. 2.13 to preserve a self-adjoint element's spectrum.\\
1.10 & $\spec(f)=f(X)$ for compact $X$ & Computes the spectrum of $f$ in the proof of Problem 2.4; illustrates the $C(X)$ model of Gelfand--Naimark.\\
1.11 & $\mathcal B(H)$ and unilateral shift & Part (a) gives a concrete model for Problem 2.1; (c,d) distinguish approximate point spectrum from spectrum, and Problem 2.10 completes the calculation of $\spec(S)$.\\
1.12 & Disc algebra & Shows that a unital, closed inclusion of Banach algebras can change the spectrum; the inclusion is not one of $C^*$-subalgebras.\\
\midrule
2.1 & $C_b(X)\hookrightarrow\mathcal B(\ell^2(X))$ & Directly verifies an isometric $*$-homomorphism; a concrete example of Thm. 2.27.\\
2.2 & Idempotents, products, unitaries & Parts (a,c) give forward directions of Prop. 2.28; (b) estimates radii of commuting products and is not needed for that proposition.\\
2.3 & Bounded product and $c_0$ direct sum & Constructs a $C^*$-algebra and its closed ideal; Problem 1.4 allows the quotient to be viewed as a Banach algebra.\\
2.4 & Continuous spectral mapping & Depends on Thm. 2.22, Prop. 2.17, and Problem 1.10; applies to functions of normal elements.\\
2.5 & $X\cong\Omega(C(X))$ & Gives uniqueness of the compact space in Cor. 2.21.\\
2.6 & Self-adjoints, unitaries, and projections & Reconstructs Prop. 2.28 through functional calculus and uses counterexamples to show why normality is needed.\\
2.7 & Unitary factor of an invertible element & Uses the positive root from Prop. 2.30 and invertibility of $a^*a$ to obtain $a=u|a|$.\\
2.8 & Continuity of $f(x_n)$ & Combines Problem 1.7, Stone--Weierstrass, and functional calculus; convergence is in norm.\\
2.9 & Monotone convergence of self-adjoint operators & Uses positive order, polarization, and Riesz; concludes strong convergence without claiming norm convergence.\\
2.10 & Spectrum of the unilateral shift & Extends Problem 1.11(d) to $\spec(S)=\overline{\mathbb D}$; $S$ is not normal.\\
\bottomrule
\end{longtable}
\normalsize

\section{Two carefully qualified bridges to Haseman}

\begin{center}
\begin{tikzpicture}[node distance=8mm]
\node[taskbox,text width=2.6cm] (q) {Problem 1.4\\quotient by an ideal};
\node[theorybox,text width=2.7cm,right=of q] (calkin) {$\K(X)$ closed\\$\mathcal B(X)/\K(X)$};
\node[researchbox,text width=2.6cm,right=of calkin] (class) {Class of $A$\\$A+\K(X)$};
\node[flowend,text width=2.7cm,right=of class] (fred) {Atkinson\\invertible $\Leftrightarrow$ Fredholm};
\draw[arr] (q)--(calkin);
\draw[arr] (calkin)--(class);
\draw[arr] (class)--(fred);
\end{tikzpicture}
\end{center}

\noindent For infinite-dimensional Banach $X$, the first arrow uses that $\K(X)$ is a closed
ideal. The last is Atkinson's theorem, not a consequence of Problem 1.4. For operators
in the Haseman research line, this formulation poses the Fredholm problem in a quotient;
see \cite[p.~2]{Karlovich2025}. As a concrete check, on $X=\ell^2$ let $P$ project onto
the first standard basis vector $e_1$, and set $A=I_X-P$. Then $A e_1=0$, but
$A+\K(X)=I_X+\K(X)$ because $P$ is compact. The calculation shows what information
the quotient retains and why Problem 1.4 prepares the formulation of Atkinson's theorem.

\begin{center}
\begin{tikzpicture}[node distance=9mm and 8mm]
\node[theorybox,text width=3.0cm] (model) {Model $L^2(\Gamma_0)$,\\$\Gamma_0=\mathbb R$ or $\mathbb T$\\$S_{\Gamma_0}^2=I$, $S_{\Gamma_0}^*=S_{\Gamma_0}$};
\node[taskbox,text width=3.0cm,right=of model] (orth) {$P_{\Gamma_0}^\pm=(I\pm S_{\Gamma_0})/2$\\orthogonal projections};
\draw[arr] (model)--(orth);
\node[researchbox,text width=3.0cm,below=of model] (plem) {Karlovich's open star\\Sokhotski--Plemelj};
\node[researchbox,text width=3.0cm,right=of plem] (trace) {$P_\pm=(I\pm S_\Gamma)/2$\\trace components};
\node[flowend,text width=3.0cm,right=of trace] (has) {$T=V_\alpha P_+ + M_GP_-$\\Haseman equation};
\draw[arr] (plem)--(trace);
\draw[arr] (trace)--(has);
\draw[mergearr,dash pattern=on 4pt off 2pt] (orth.south)--(trace.north);
\end{tikzpicture}
\end{center}

\noindent The dashed arrow compares only the form of the formulas. In the model of
Bastos--Fernandes--Karlovich \cite[§2, p.~22]{BFK2006}, $P_{\Gamma_0}^\pm$ are
orthogonal projections. On the open star of Karlovich 2017, boundary reduction to $T$
follows from Sokhotski--Plemelj \cite[pp.~463--464]{Karlovich2017}; according to
Theorem 5.3 of that source, $S_\Gamma^2=I$ does not hold there in general. Thus the
operators $P_\pm$ in that equation are not called idempotents or projections here.
For the check $\alpha=\mathrm{id}$ and $G=1$, the Haseman formula gives
$T=P_++P_-=I$ by cancellation of the two $S_\Gamma$ terms. This calculation uses
the Plemelj formula without requiring $S_\Gamma^2=I$.

\begin{reflection}
To check an arrow, locate the line of the proof that uses its source. For example, identify
where Problem 1.9 enters the reality of the spectrum and where Problem 1.10 enters Problem
2.4. Then verify in the $L^2(\Gamma_0)$ model how $S_{\Gamma_0}^2=I$ gives
idempotents and how $S_{\Gamma_0}^*=S_{\Gamma_0}$ makes them orthogonal projections.
Compare those assumptions with the assumptions for the open star.
\end{reflection}
